\section{Design Participativo e Validação Territorial}

Este módulo incorpora usuários finais na co-criação das regras e interfaces do sistema, avaliando usabilidade, apropriação e adoção continuada.

\subsection{Fundamentação Metodológica}

O Design Participativo (Participatory Design -- PD) constitui abordagem centrada no usuário que promove envolvimento ativo dos stakeholders ao longo de todo o ciclo de desenvolvimento de tecnologias, desde a definição de requisitos até a validação \cite{Schuler1993,Simonsen2012}. Para sistemas de apoio à decisão em contextos comunitários, o PD é crítico para garantir: (i) aderência dos requisitos às necessidades reais; (ii) apropriação psicológica da tecnologia; (iii) capacitação para uso autônomo; e (iv) sustentabilidade do sistema após a saída da equipe de pesquisa.

\subsection{Delineamento}

O processo participativo será estruturado em cinco fases:

\subsubsection{Fase 1: Diagnóstico Participativo de Necessidades}

Realização de grupos focais (n~=~3--4, com 8--12 participantes cada) nas comunidades quilombolas do Sul de Sergipe para identificar necessidades de informação, desafios na valoração e comercialização de produtos, e expectativas em relação a ferramentas de apoio.

\textbf{Técnicas:} mapeamento de stakeholders, análise SWOT participativa, construção de personas (perfis típicos de usuários) e jornadas de usuário (user journeys).

\subsubsection{Fase 2: Co-Criação de Requisitos e Regras}

Oficinas de co-design (n~=~2--3, duração 4--6h cada) onde os participantes, juntamente com a equipe de pesquisa, definirão:

\begin{itemize}
\item \textbf{Priorização de variáveis:} quais dimensões de valor são mais relevantes para inclusão no IVB (complementando o painel Delphi com perspectiva comunitária).
\item \textbf{Calibração de limiares:} ajuste dos intervalos de funções de pertinência fuzzy ("o que é 'Alta Autenticidade' na prática?").
\item \textbf{Refinamento de regras:} validação e ajuste de regras de inferência, verificando coerência com conhecimento local.
\item \textbf{Design de interface:} prototipagem em papel (low-fidelity prototypes) das telas do sistema, definindo fluxo de navegação, terminologia acessível e visualizações de saída (gráficos, tabelas).
\end{itemize}

\textbf{Ferramentas:} uso de cartões, diagramas impressos, sketches e post-its para facilitar ideação colaborativa. Sessões gravadas em áudio/vídeo (com consentimento) para análise posterior.

\subsubsection{Fase 3: Prototipagem Iterativa e Testes de Usabilidade}

Desenvolvimento de protótipo funcional de média fidelidade (implementado em Python com interface Streamlit ou similar) incorporando os requisitos co-criados. O protótipo será submetido a testes de usabilidade em laboratório e em campo (nas comunidades).

\textbf{Protocolo de Teste:}
\begin{enumerate}
\item \textbf{Tarefas dirigidas:} os usuários (n~=~10--15) executarão tarefas típicas (ex.: "Calcule o IVB para o produto X", "Compare dois produtos", "Gere relatório para apresentar a potenciais compradores").
\item \textbf{Métricas quantitativas:} taxa de sucesso nas tarefas, tempo de conclusão, número de erros, necessidade de assistência.
\item \textbf{Métricas qualitativas:} entrevistas pós-teste explorando dificuldades, sugestões de melhoria e satisfação geral (System Usability Scale -- SUS \cite{Brooke1996}).
\item \textbf{Análise think-aloud:} usuários verbalizarão pensamentos durante execução de tarefas, permitindo identificar pontos de confusão.
\end{enumerate}

Ciclos iterativos de redesign e teste serão conduzidos até alcançar taxa de sucesso $\geq$ 85\% e SUS $\geq$ 70 (limiar de usabilidade aceitável).

\subsubsection{Fase 4: Implantação Piloto e Capacitação}

Implantação do sistema finalizado em 2--3 comunidades quilombolas, acompanhada de programa de capacitação (oficinas de 8--12h) cobrindo:

\begin{itemize}
\item Fundamentos de lógica fuzzy e valoração de ativos (conceitual, não técnico).
\item Operação do sistema: entrada de dados, interpretação de resultados, geração de relatórios.
\item Estratégias de uso: integração do IVB em negociações comerciais, comunicação com consumidores, suporte a processos de certificação.
\end{itemize}

\textbf{Materiais:} manual do usuário ilustrado, vídeos tutoriais curtos (3--5 min cada), e linha de suporte técnico (WhatsApp/telefone).

\subsubsection{Fase 5: Avaliação de Adoção e Apropriação}

Após 6--12 meses de uso, será conduzida avaliação de adoção e apropriação mediante:

\begin{itemize}
\item \textbf{Métricas de uso:} frequência de acesso ao sistema, número de produtos valorados, funcionalidades mais utilizadas (logs do sistema).
\item \textbf{Apropriação psicológica:} aplicação da escala de apropriação tecnológica \cite{Carroll2001}, mensurando sentimento de propriedade, adaptação criativa e evangelização (recomendação a outras comunidades).
\item \textbf{Impacto percebido:} entrevistas e questionários explorando se o sistema contribuiu para: melhoria de preços, acesso a novos mercados, fortalecimento de argumentação em negociações, empoderamento na gestão de ativos.
\item \textbf{Sustentabilidade:} avaliação de capacidade de uso autônomo (sem suporte contínuo da equipe de pesquisa) e manutenção do sistema.
\end{itemize}

\subsection{Indicadores de Validação}

\begin{itemize}
\item \textbf{Usabilidade:} SUS $\geq$ 70 (aceitável) ou $\geq$ 80 (bom).
\item \textbf{Taxa de adoção:} $\geq$ 60\% dos participantes da capacitação utilizando o sistema regularmente (ao menos 1x/mês) após 6 meses.
\item \textbf{Apropriação:} $\geq$ 50\% dos usuários reportando adaptações criativas do sistema (uso não previsto inicialmente) ou evangelização ativa.
\item \textbf{Impacto percebido:} $\geq$ 70\% dos usuários reportando contribuição do sistema para algum resultado concreto (melhoria de preço, acesso a mercado, fortalecimento de argumentação).
\end{itemize}

A validação territorial bem-sucedida constituirá evidência de que o sistema não apenas produz saídas tecnicamente corretas, mas é efetivamente utilizável, apropriável e útil no contexto real de governança e comercialização de ativos tradicionais.
